\begin{tikzpicture}[
    tensor/.style={rectangle, draw, fill=white, text width=2em, text centered, minimum height=2em},
    % % --- NUEVOS ESTILOS PARA ETIQUETAS ---
    lbl_fwd/.style={above, pos=0.5, font=\footnotesize}, % 'pos=0.5' centra la etiqueta
    lbl_bwd/.style={below, pos=0.5, font=\footnotesize},
    lbl_left/.style={left, pos=0.5, font=\footnotesize},
    lbl_right/.style={right, pos=0.5, font=\footnotesize},
    layer_function/.style={rectangle, draw, fill=white, text width=4em, text centered, rounded corners, minimum height=2em},
    operation/.style={circle, draw, fill=white, minimum size=0.1em, font=\tiny},
    arrow/.style={-latex, thick}
]

% Distancia global de nodos.
\tikzset{node distance=1.75cm and 1.5cm}

%%%% FORWARD %%%%
% draw linear block z()
\node[layer_function] (z) at (0,0) {$z()$};

% Input activation
\draw[arrow] ($(z.west)+(-2.5,0)$) -- (z.west) node[lbl_fwd] {$a_{m-1}^{(q)}$};

% --- Cadena de activación con cuantización ---
\node[layer_function, right=of z] (phi) {$\phi()$};
\node[layer_function, right=of phi] (qa) {$q_\phi()$};
\draw[arrow] (z.east) -- (phi.west);
\draw[arrow] (phi.east) -- (qa.west);
\draw[arrow] (qa.east) -- node[lbl_fwd, pos=0.75] {$a_{m}^{(q)}$} ($(qa.east)+(2.5,0)$);

% --- Cadena de parámetros con cuantización ---
\node[layer_function, below=0.75cm of z] (qtheta) {$q_{\theta}()$};
\node[tensor, below=0.75cm of qtheta] (theta) {$\theta$};
\draw[arrow] (qtheta) -- (z);
\draw[arrow] (theta) -- (qtheta);


% % --- DEFINE PARAMETER UPDATE CHAIN FIRST --- % %
\def\updateopspacing{0.5cm}
\node[operation, below=\updateopspacing of theta] (mult_eta) {$\times$};
\node[operation, below=\updateopspacing of mult_eta] (negative) {$-$};
\node[left=0.5cm of mult_eta] (eta) {$\eta$};


%%%% BACKWARD %%%%

% --- Cadena vertical de gradientes ---
\node[layer_function, below=0.75cm of negative] (dqtheta) {$\nabla_{q_{\theta}}$};
\node[layer_function, below=0.75cm of dqtheta] (dz) {$\nabla_z$};

% --- Cadena horizontal de gradientes ---
% Usamos distancias manuales 'right=Xcm' para compactar esta parte
\node[layer_function, right=1.25cm of dz] (dphi) {$\nabla_\phi$};
\node[layer_function, right=1.25cm of dphi] (dqa) {$\nabla_{q_\phi}$};

% --- Posición de 'mult' y su flecha ---
\node[operation, right=0.75cm of dqa] (mult) {$\times$};
% Input grad (dL/da_m)
\draw[arrow] ($(mult.east)+(2.0,0)$) -- (mult.east) node[lbl_bwd] {$\frac{\partial \mathcal{L}}{\partial a_{m}^{(q)}}$};
\draw[arrow] (mult.west) -- (dqa.east);

% Multiplicadores en cadena
\node[operation] (mult3) at ($(dphi)!.5!(dqa)$) {$\times$};
\draw[arrow] (dqa.west) -- (mult3.east);
\draw[arrow] (mult3.west) -- (dphi.east);

\node[operation] (mult2) at ($(dz)!.5!(dphi)$) {$\times$};
\draw[arrow] (dphi.west) -- (mult2.east);
\draw[arrow] (mult2.west) -- (dz.east);

% Output grad (dL/da_{m-1})
\draw[arrow] (dz.west) -- node[lbl_bwd] {$\frac{\partial \mathcal{L}}{\partial a_{m-1}^{(q)}}$} ($(dz.west)+(-2.5,0)$);


% % --- CONNECT ALL THE NODES --- % %
\draw[arrow] (eta.east) -- (mult_eta.west);
\draw[arrow] (mult_eta.north) -- (theta.south);
\draw[arrow] (negative.north) -- (mult_eta.south);

% --- Conexiones verticales de gradiente ---
\draw[arrow] (dqtheta.north) -- node[lbl_left] (dtheta) {$\frac{\partial \mathcal{L}}{\partial \theta}$} (negative.south);
\draw[arrow] (dz.north) -- node[lbl_left] (dq) {$\frac{\partial \mathcal{L}}{\partial q_{\theta}}$} (dqtheta.south);


% % --- BACKGROUND LAYER BOX --- % %
\begin{pgfonlayer}{background}
    \node [
        draw,
        dotted,
        thick,
        orange,
        rounded corners,
        fit=(z) (phi) (qa) (theta) (qtheta) (mult_eta) (negative) (dz) (dphi) (dqa) (mult) (mult2) (mult3) (dtheta) (dq) (eta),
        inner sep=0.35cm
    ] (layer_box) {};

    % --- MODIFICADO: La línea cruza exactamente por el centro de THETA ---
    % Al estar en la capa 'background', el nodo theta (que es blanco) tapará la línea,
    % dando el efecto de que pasa por detrás.
    \draw [gray, dashed] ($(theta)+(-2.5,0)$) -- ($(theta)+(10,0)$);

\end{pgfonlayer}

% Label del layer
\node[font=\footnotesize, orange, above right, xshift=2mm] at (layer_box.north east) {$L_m$};

\end{tikzpicture}
